\section{Síntesis y ampliación}

\subsection{Repaso de los puntos clave}

Esta clase parte del modelo de lenguaje n-gram más elemental y, paso a paso, revela los principios de funcionamiento y las prácticas de ingeniería de los LLM modernos. Los puntos clave pueden resumirse en los siguientes:

\begin{enumerate}
    \item \textbf{Todo comienza con next-word prediction}: desde los modelos de lenguaje Bayesian de la década de 1980 hasta Gemini en 2025, la tarea central no ha cambiado: predecir el siguiente token. Todo comportamiento que parece «inteligente» es un resultado emergente de este objetivo simple.

    \item \textbf{Scale produce un cambio cualitativo}: el crecimiento exponencial del número de parámetros (de millones a billones) y de la ventana de contexto (de 4 palabras a 2 millones de tokens) dio lugar a capacidades emergentes como few-shot learning.

    \item \textbf{Un LLM es una consulta difusa sobre los datos de entrenamiento}: la salida del modelo es, en esencia, la reproducción de patrones estadísticos de los datos de entrenamiento. La eficacia del Role Prompting (que «MIT mathematician» eleve la tasa de aciertos en problemas matemáticos) confirma esto: solo estamos condicionando distintas regiones de la distribución de salida.

    \item \textbf{El Prompt Engineering es una manipulación del espacio de probabilidad}: Role Prompting condiciona la distribución; Chain-of-Thought amplía la superficie de error; las indicaciones de formato guían al modelo hacia un patrón de datos específico.

    \item \textbf{El Post-training desplaza al modelo de lenguaje}: Instruction Tuning, RLHF y Constitutional AI son medios para moverse entre «modelos de lenguaje válidos»; en esencia, actualizan los pesos para hacer shift a la distribución de probabilidad.

    \item \textbf{Los LLM tienen limitaciones graves}: alucinaciones, sesgos, Prompt Injection, incumplimiento de reglas y aceptación sin cuestionar premisas erróneas. El uso de LLM debe ir acompañado de una evaluación rigurosa y de mecanismos de seguridad externos.

    \item \textbf{Agent = razonamiento + uso de herramientas}: ReAct combina razonamiento y acción, y Toolformer enseña al modelo a usar herramientas externas. Estos dos artículos sentaron las bases teóricas de los sistemas Agent modernos.
\end{enumerate}

\subsection{Lecturas adicionales}

\begin{itemize}
    \item \textbf{Artículo de GPT-3}: Brown et al., «Language Models are Few-Shot Learners», NeurIPS 2020 —— sentó las bases del few-shot prompting
    \item \textbf{Artículo de Chinchilla}: Hoffmann et al., «Training Compute-Optimal Large Language Models», 2022 —— la proporción óptima entre datos de entrenamiento y número de parámetros
    \item \textbf{Chain-of-Thought}: Wei et al., «Chain-of-Thought Prompting Elicits Reasoning in Large Language Models», NeurIPS 2022
    \item \textbf{Zero-shot CoT}: Kojima et al., «Large Language Models are Zero-Shot Reasoners», NeurIPS 2022 —— el descubrimiento de «Let's think step by step»
    \item \textbf{LoRA}: Hu et al., «LoRA: Low-Rank Adaptation of Large Language Models», ICLR 2022
    \item \textbf{RLHF}: Ouyang et al., «Training language models to follow instructions with human feedback», NeurIPS 2022
    \item \textbf{Constitutional AI}: Bai et al., «Constitutional AI: Harmlessness from AI Feedback», Anthropic 2022
    \item \textbf{ReAct}: Yao et al., «ReAct: Synergizing Reasoning and Acting in Language Models», ICLR 2023
    \item \textbf{Toolformer}: Schick et al., «Toolformer: Language Models Can Teach Themselves to Use Tools», NeurIPS 2023
    \item \textbf{LearnPrompting.org}: recurso de aprendizaje sobre Prompt Engineering recomendado por el docente
\end{itemize}

%% --- Fin del contenido --- %%

\end{document}
